%% notes-tex/019-simb-multimodal-expression/sections/5-joint-plan.tex
%% SOURCE: the five reviewer reports and the coordinator's addendum of 2026-09-27 under
%% review/2026-09-27-joint-review/ (01 data ceilings, 02 trunk and collapse, 03 joint
%% evidence, 04 statistics and design, 05 literature and designs, 05a the mRNA-protein
%% survey, launch_census.json), each of which names the W&B run, results file or mirror
%% paper behind every number; the v16 readouts results/v16_joint_readout.json and
%% v16_joint_expr_readout.json; the weekly note 2026.09.27.

\section{Joint proteome and expression training: what is established, what the data permits, and the rounds that follow\secstatus{todo}}
\label{sec:joint-plan}

The question, as posed on 2026-09-25: can training one trunk on the knockout proteome and
the knockout transcriptome be shown to help each label's own prediction, with a statement
that survives review, inside two weeks of compute on the IGB \file{gpu} and \file{cabbi}
partitions. Five independent reviews of the campaign's runs, code and the literature were
run on 2026-09-27 and are filed under \file{review/2026-09-27-joint-review/}; this section
is their synthesis. Every number is measured unless it is marked as a hypothesis.

\pfstep{What v16 measured, read again} On the proteome head the joint arm against the
proteome-only arm is $-0.008$ at the fixed window ($t = -1.4$, three of six pairs positive),
$-0.016$ on the rolling maximum (none of six), $+0.005$ on the held-out test set at the
best-validation checkpoint (four of six), and $-0.095$ on the training fit at epoch 499
(none of six, $t = -15$). The joint arm's proteome head launches later than the
proteome-only head in every pair (epochs 95 to 144 against 44 to 61). The proteome-only arm
at 500 epochs reproduces v14's reference at 1,353 to 1,960 epochs to within $0.005$ on all
six split-and-seed pairs, so the proteome head is budget-saturated at 500 and a longer budget
cannot move that contrast. On the expression head the published $+0.060$ is a comparison
against a control that never trained (\cref{sec:split-rounds}); on the two clean pairs it is
$+0.007$ on validation and $-0.067$ on test. A properly trained expression-only model from
v13, cut at epoch 500, scores in the same band as the joint arm's expression head at 500
(rolling maxima $0.064$ to $0.124$ against $0.082$ to $0.118$), so the joint head sits on the
single-head curve rather than ahead of it. One result the readouts do not report: with ten
proteins revealed through the masked-conditioning channel the joint arm's proteome head beats
the proteome-only arm by $+0.090$ in six of six pairs ($t = 20$), and with 1,000 revealed it
loses by $0.036$ in six of six. That is a conditioning result and it rides on the objective
that is being dropped. Nothing measured supports the target claim, and the proteome side is
directionally against it.
\src{review/2026-09-27-joint-review/03_joint_evidence.md, 02_trunk_collapse.md}

\pfstep{Why the single-head control died, and how to keep it alive} After per-gene
standardization, predicting each gene's marginal quantile is an exact strain-blind stationary
point of the pinball loss, reachable from the strain-invariant gene state alone, and it sits
only about $3\,\%$ above the working model in loss ($0.239$ against $0.231$ on one card).
Whether a head leaves that point is a signal-to-noise race in the number of labeled rows per
optimizer step, and the pinball gradient carries no residual magnitude to help it. Across 95
launched expression runs on the undiluted stores the launch epoch never exceeded 314 and its
90th percentile is at most 115; the v16 expression-only arm, at eleven labeled rows per batch,
launched twice, at epochs 138 and 195, and four times never. A resumed checkpoint does not
escape: the continuation of \texttt{hyasw3nx} is still at zero at epoch 555. A detector on the
logged prediction spread (below $0.05$ by epoch 200) catches every never-launched run with no
false positive in 110 healthy runs, and a run that has trained for a thousand epochs can fall
back into the same state (\texttt{p204pwqb}), so the detector runs for the whole run. Not
implicated on the evidence: gradient clipping (never binds), the scheduler (none exists), the
zero-gated pathways (every gate in these rounds is forced open), co-residency, DDP, the
missing-value masking. Warmup and the ReZero residual have never been measured; the only four
runs that carried them died inside the pre-launch plateau in July.
\src{review/2026-09-27-joint-review/02_trunk_collapse.md, launch_census.json}

\pfstep{What the data permits a shared trunk to transfer} On the 1,349 strains carrying both
labels, a plain ridge from the observed proteome predicts held-out expression at per-feature
Pearson $0.226$, against $0.101$ from the deleted gene's ProtT5 embedding; the reverse reads
$0.218$ against $0.056$. The two channels barely overlap: adding the genotype to the observed
proteome buys $+0.009$. A trunk that sees only the genotype at scoring time cannot collect that
channel. What it can collect is a shared genotype representation, and the linear measurement of
that channel is $+0.005$ to $+0.006$ on the proteome head (the expression task's top-128 and
top-256 ridge directions beat the proteome's own) and at or below zero on the expression head
(the proteome's top-8 directions deliver $0.044$ against $0.110$ from expression's own, with a
random floor of $0.030$). The two tasks share three to thirteen times more genotype directions
than chance and keep 63 to 90 percent of each subspace private. Per-deletion agreement between
the two modalities, $0.036$, equals the agreement between duplicate Messner strains of the
same deletion, $0.03$, so it carries no information; the per-gene agreement, $0.075$
(disattenuated $0.15$ to $0.21$), and the gene-pair co-variation agreement of $0.31$ are the
transferable structure. The proteome is the noisier label: 28 percent of the 1,850 proteins
have knockout variance at or below the replicate variance, and per-feature standardization
gives those columns full weight in both the loss and the metric. Detecting $+0.005$ at 80
percent power needs 61 to 88 pairs at the measured paired sd of $0.0167$; no fourteen-day
design reaches that.
\src{review/2026-09-27-joint-review/01_data_ceilings.md, 05a_literature_helper_mrna_protein.md}

\pfstep{What the field has shown} The one same-model joint-against-single test in the
multimodal literature, totalVI against scVI on held-out RNA from the same cells, is null. The
one measured pairing ladder (Ryu et al., paired RNA and protein in the same T cells) values the
per-sample pairing at about $0.028$ Pearson even with a perfect pairing. The sign of a
multi-task result flips on the capacity control (Standley et al.: worse at equal per-task
capacity, better at equal total compute), so the control has to be declared before the round.
The field's one consistent positive for perturbation prediction is a linear model with a
perturbation embedding pretrained on other perturbation data (Ahlmann-Eltze et al.), a transfer
design. The Messner-against-Kemmeren per-deletion comparison has never been published; the
closest is a strain-level $R = 0.51$ in fission yeast on 94 paired knockouts.
\src{review/2026-09-27-joint-review/05_literature_designs.md}

\pfstep{The statistic that can carry a conjunction} A claim that joint training helps both
heads is an intersection-union test: reject only when both one-sided tests reject, which holds
level $\alpha$ with no multiplicity correction. Its price is power: if the true proteome
effect is exactly zero, a superiority-and-superiority rule has five percent power at any
sample size. The defensible form is superiority on the expression head and non-inferiority on
the proteome head at a margin fixed before launch; $0.015$ sits above the between-partition sd
of the proteome score and below the single-partition bootstrap sd of $0.026$. At the v16 point
estimate of $-0.008$, non-inferiority at margin $0.010$ needs 275 pairs, at $0.015$ needs 24,
at $0.020$ needs 10. Two more measured facts shape the design. The arm effect clusters by
partition (intraclass correlation $0.50$ to $0.72$ on two contrasts), so the paired $t$ over
twelve split-by-seed cells understates its standard error by 1.5 to 1.7 times, and a
partition-level sign-flip test has an attainable minimum $p$ of $2/2^P$, which four partitions
can never bring under $0.05$; at a fixed run count more partitions always beat more seeds, and
36 runs as twelve partitions with one seed each give a standard error of $0.0070$ against
$0.0098$ for four by three. And the committed readout applies its plateau rule to the rolling
maximum rather than to the registered window, which is the one line that decides whether v17's
propagation contrast reads $+0.012$ ($p = 0.23$) or $+0.020$ on eleven pairs ($p = 0.004$);
it is reported here as found, not changed, because changing a registered rule after the fact
is the error the rule exists to prevent.
\src{review/2026-09-27-joint-review/04_statistics_design.md}

\pfstep{What is provable in two weeks, and what is not} Superiority of the joint trunk on
both heads is not, unless the effect is larger than anything this trunk has shown: the linear
transfer bound is $+0.006$ on one head and zero on the other, the minimum detectable effect at
twelve pairs is $0.013$ to $0.019$, and the external priors are against it. Three statements
are within reach. First, that joint training costs neither head more than a pre-registered
margin while one model serves both labels, which twelve to eighteen clean pairs settle.
Second, that a model given a strain's proteome predicts its transcriptome far better than a
model given the genotype alone, and the reverse, an effect of $+0.11$ to $+0.16$ that six
pairs establish beyond doubt; this is a conditioning claim about a multimodal model, not a
shared-trunk synergy claim, and it has to be labeled as such. Third, a clean answer to the
synergy question at a resolution of $0.013$, with a permutation control that says whether any
gain uses the paired labels or only their marginal structure, and with the single-head
controls kept alive; a null there is a measurement with a stated resolution, not a failure.

\pfstep{The rounds} Zero-GPU work first, one to two days, all of it prerequisite: a launch
guard callback that kills a run whose prediction spread has not passed $0.05$ by epoch 200 and
relaunches it from a new initialization inside the same cell, recording the relaunch; a test
pass per head at that head's own best-validation checkpoint (the joint arm's expression test
number is currently read at the proteome's checkpoint, $0.032$ low on average); a
label-permutation option for one head across training strains; gradient-cosine logging between
the two heads on the shared trunk; the gene-disjoint expression subset and a
reliability-weighted proteome Pearson as logged metrics; and the pre-registration written into
the round's config header. Then two rounds, launchable together.

\begin{itemize}
\item \textbf{v19, the deconfounded joint round.} Store \file{fig3_proteome} restricted to the
  strains carrying both labels (about 1,100 training rows per partition, 35 steps per epoch,
  three times cheaper than the union), so every arm sees identical rows and steps. Four arms on
  one card, derived from one configuration by loss weights alone: proteome head only,
  expression head only, both heads, and both heads with the expression labels permuted across
  training strains. Every head in the unmasked form the joint arm's auxiliary head uses; the
  masked objective off in every arm, by the decision of 2026-09-25, with its measured cost of
  $0.012$ borne equally. Twelve split seeds, one initialization seed each, 1,200 epochs,
  per-head windows declared from v14 and v16 (proteome epochs 200 to 400, expression 1,000 to
  1,200). Primary test: expression superiority and proteome non-inferiority at margin $0.015$,
  one-sided $0.05$ each, paired over the twelve partitions; secondary: the interaction with the
  permuted arm on both heads, the partition-level sign-flip $p$, intent-to-treat beside the
  guarded analysis, the test read per head. Twelve tasks of four runs at about 32 hours each
  (v13 ran four per A40 at 40 steps per epoch, 1,200 epochs in about 32 hours): on three
  \file{gpu} cards under the \texttt{\%3} rule plus two \file{cabbi} cards, three waves, about
  five days, two more for relaunches, every job inside the five-day limit.
\item \textbf{v20, the conditioned cross-modal round.} Same store. Expression head from the
  genotype alone against the same head with the strain's observed proteome revealed through the
  existing conditioning channel and scored with it revealed; the mirror pair predicting the
  proteome with expression revealed. The first direction needs no training: the six v16 joint
  runs already trained with the proteome revealed on three of four steps and validation already
  sweeps the reveal sizes, but the trainer logs the sweep for the masked head only, so one
  change to \file{\_cache\_masked\_metric} and a re-evaluation of the six saved checkpoints at
  zero against 1,000 revealed proteins on the expression head, with another strain's proteome as
  the permuted control, gives six within-run pairs in under a GPU-day. The mirror direction is a
  training round: three split seeds, two initialization seeds, two arms, 600 epochs, three per
  card, four tasks of about 20 hours on \file{cabbi}, about two days. Expected effects $+0.11$
  and $+0.16$ against a minimum detectable $0.019$.
\end{itemize}

Two decisions are the author's. Whether v19 runs on the both-label intersection, which makes
the contrast clean and cheap, or on the union, where the joint arm also gains 2.5 times more
strains and the claim becomes ``joint training on the union helps'' with the data-quantity term
inside it; a fifth arm on a subset of partitions can price that term. And what to do with the
v16 continuation (IGB 2410399, split 0 at epoch about 600, four tasks pending): it extends the
confounded control and cannot move the proteome contrast, so canceling the pending tasks frees
the \file{gpu} cards for v19 as soon as the prerequisites land.
\src{review/2026-09-27-joint-review/04_statistics_design.md, 03_joint_evidence.md, 01_data_ceilings.md, 05_literature_designs.md}
